\section{Experiment Results} \label{sec:exp}
In this section, we will demonstrate the strong empirical performance of our proposed meta-model-based uncertainty learning method: first, We introduce the UQ applications; then we describe the experiment settings; next, we present the main results of the three aforementioned uncertainty quantification applications; and finally, we discuss our takeaways. More experiment results and implementation details are given in the Appendix~\ref{setup-appendix} and Appendix~\ref{results-appendix}.


\subsection{Uncertainty Quantification Applications} \label{sec:app}
We primarily focus on three applications that can be tackled using our proposed meta-model approach: (1) \textbf{Out of domain data detection.}
Given a base-model $\vh$ trained using data sampled from the distribution $P^B_{Z}$. We use the same base-model training set to train the meta-model, i.e., $\mathcal{D}_{B} = \mathcal{D}_{M}$. During testing, there exists some unobserved out-of-domain data from another distribution $P^{ood}_{Z}$. The meta-model is expected to identify the out-of-distribution input sample based on epistemic uncertainties.  (2) \textbf{misclassification Detection.}
Instead of detecting whether a testing sample is out of domain, the goal here is to identify the failure or success of the meta-model prediction at test time using total uncertainties. (3) \textbf{Trustworthy Transfer Learning.}
%The aforementioned applications assume the data used for training the base-model and the meta-model are exactly same, i.e., $\mathcal{D}_{B} = \mathcal{D}_{M}$, or $P^B_{Z} = P^B_{M}$ in general. Instead, 
In transfer learning, there exists a pretrained model trained using source task data $\mathcal{D}_{s}$ sampled from source distribution $P^s_{Z}$, and the goal is to adapt the source model to a target task using target data $\mathcal{D}_{t}$ sampled from target distribution $P^t_{Z}$. Most existing transfer learning approaches only focus on improving the prediction performance of the transferred model, but ignore its UQ performance on the target task. Our meta-model method can be utilized to address this problem, i.e, given pretrained source model $\vh^{s} \circ \Phi^{s}$, the meta-model can be efficiently trained using target domain data by $\vg^t = \argmin_{\vg} \gL_E(\vg \circ \Phi^s,\gD_t)$.


\begin{table*}[t]
  \begin{center}
  \scriptsize
  \captionsetup{font=small}
  \caption{\textbf{OOD Detection AUROC score.} MI, Dent, and Prec stand for different epistemic uncertainty metrics, i.e., Mutual Information, Differential Entropy, and precision. Settings stand for post-hoc or traditional, i.e., training the entire model from scratch. Additional Data stands for if using additional training data or not.}
%   MI, Dent, and Prec stand for different epistemic uncertainty metrics, i.e., Mutual Information, Differential Entropy, and precision.
  \vspace{-2em}
  \begin{tabular}{lllllllll}
    \\
    \toprule
    \textbf{ID Data}\ \&\ \textbf{Model} & \textbf{Methods} & \textbf{Settings} & \textbf{Additional Data} & \textbf{Omniglot}& \textbf{FMNIST} & \textbf{KMNIST}& \textbf{CIFAR10} &\textbf{Corrupted} \\
    \hline
    MNIST & Base Model(Ent) & Traditional & No & 98.9$\pm$0.5 & 97.8$\pm$0.8 & 95.8$\pm$0.8 & 99.4$\pm$0.2 & 99.5$\pm$0.3  \\
    LeNet & Base Model(MaxP) & Traditional & No & 98.7$\pm$0.6 & 97.6$\pm$0.8 & 95.6$\pm$0.8 & 99.3$\pm$0.2 & 99.4$\pm$0.4  \\
              & Whitebox & Post-hoc & Yes & 98.5$\pm$0.3 & 97.7$\pm$0.6 & 96.0$\pm$0.2 & 99.5$\pm$0.1 & 99.5$\pm$0.1  \\
              & LULA & Post-hoc & Yes & 99.8$\pm$0.0 & 99.4$\pm$0.0 & \textcolor{violet}{\textbf{99.3}$\pm$0.1} & 99.9$\pm$0.0 & 99.6$\pm$0.1  \\
    \hline
              & \textbf{Ours(Ent)} & Post-hoc & No & 99.7$\pm$0.1 & 99.5$\pm$0.2 & 98.2$\pm$0.3 &  \textcolor{violet}{\textbf{100.0}$\pm$0.0} & \textcolor{violet}{\textbf{100.0}$\pm$0.0} \\
              & \textbf{Ours(MaxP)} & Post-hoc & No & 99.3$\pm$0.2 & 99.3$\pm$0.2 & 98.0$\pm$0.2 &  \textcolor{violet}{\textbf{100.0}$\pm$0.0} & \textcolor{violet}{\textbf{100.0}$\pm$0.0} \\
              & \textbf{Ours(MI)} & Post-hoc & No & \textcolor{violet}{\textbf{99.9}$\pm$0.0} & \textcolor{violet}{\textbf{99.6}$\pm$0.2} &97.7$\pm$0.4 &  \textcolor{violet}{\textbf{100.0}$\pm$0.0} & \textcolor{violet}{\textbf{100.0}$\pm$0.0} \\
              & \textbf{Ours(Dent)} & Post-hoc & No & 99.8$\pm$0.0 & 99.5$\pm$0.2 & 97.6$\pm$0.4 & \textcolor{violet}{\textbf{100.0}$\pm$0.0} & \textcolor{violet}{\textbf{100.0}$\pm$0.0} \\
              & \textbf{Ours(Prec)} & Post-hoc & No & \textcolor{violet}{\textbf{99.9}$\pm$0.0} & \textcolor{violet}{\textbf{99.5}$\pm$0.2} &97.7$\pm$0.5 &  \textcolor{violet}{\textbf{100.0}$\pm$0.0} & \textcolor{violet}{\textbf{100.0}$\pm$0.0} \\
    \toprule
    \textbf{ID Data}\ \&\ \textbf{Model} & \textbf{Methods} & \textbf{Settings} & \textbf{Additional Data} & \textbf{SVHN}& \textbf{FMNIST} & \textbf{LSUN}& \textbf{TinyImageNet} &\textbf{Corrupted} \\
    \hline
    CIFAR10 & Base Model(Ent) & Traditional & No & 86.4$\pm$4.6 & 90.8$\pm$1.3 & 89.0$\pm$0.5 & 87.5$\pm$1.1 & 85.9$\pm$8.2  \\
    VGG16   & Base Model(MaxP) & Traditional & No & 86.3$\pm$4.4 & 90.4$\pm$1.2 & 88.7$\pm$0.5 & 87.3$\pm$1.1 & 85.7$\pm$8.1   \\
              & Whitebox & Post-hoc & Yes & 96.9$\pm$0.9 & 95.2$\pm$1.2 & 89.3$\pm$2.2 & 88.9$\pm$2.5 & 96.4$\pm$1.0   \\
              & LULA & Post-hoc & Yes & 97.1$\pm$1.7 & 94.3$\pm$0.0 & 92.8$\pm$0.1 & 90.0$\pm$0.0 & 97.7$\pm$2.0  \\
    \hline
              & \textbf{Ours(Ent)} & Post-hoc & No & 96.3$\pm$3.0 &89.0$\pm$5.2 & 89.6$\pm$3.4 &  89.4$\pm$3.5& 95.9$\pm$4.3 \\
              & \textbf{Ours(MaxP)} & Post-hoc & No & 95.6$\pm$3.6 & 87.8$\pm$4.4 &89.1$\pm$2.4 &  88.2$\pm$2.6 & 94.0$\pm$7.3 \\
              & \textbf{Ours(MI)} & Post-hoc & No & \textcolor{violet}{\textbf{100.0$\pm$0.0}} & \textcolor{violet}{\textbf{98.8$\pm$0.5}} & 95.2$\pm$0.9 &  \textcolor{violet}{\textbf{98.1$\pm$0.3}} & \textcolor{violet}{\textbf{100.0$\pm$0.0}} \\
              & \textbf{Ours(Dent)} & Post-hoc & No & \textcolor{violet}{\textbf{100.0$\pm$0.0}} & 98.4$\pm$0.9 & \textcolor{violet}{\textbf{95.7}$\pm$0.8} &  97.7$\pm$0.5 & \textcolor{violet}{\textbf{100.0$\pm$0.0}} \\
               & \textbf{Ours(Prec)} & Post-hoc & No & \textcolor{violet}{\textbf{100.0$\pm$0.0}} & \textcolor{violet}{\textbf{98.8$\pm$0.5}} & 95.1$\pm$0.5 &  \textcolor{violet}{\textbf{98.1$\pm$0.3}} & \textcolor{violet}{\textbf{100.0$\pm$0.0}} \\
    \toprule
    CIFAR100 & Base Model(Ent) & Traditional & No & 76.2$\pm$5.2 & 77.8$\pm$2.4 & 80.1$\pm$0.5 & 79.7$\pm$0.3 & 65.8$\pm$11.4  \\
    WideResNet   & Base Model(MaxP) & Traditional & No & 73.9$\pm$4.3 & 76.4$\pm$2.3 & 78.7$\pm$0.5 & 78.0$\pm$0.2 & 63.8$\pm$10.4   \\
              & Whitebox & Post-hoc & Yes & 89.0$\pm$0.7 & 82.4$\pm$1.1 & 80.5$\pm$0.7 & 79.0$\pm$1.1 & 83.1$\pm$1.6   \\
              & LULA & Post-hoc & Yes & 84.2$\pm$1.0 & 83.2$\pm$0.1 & 79.6$\pm$0.3 & 78.5$\pm$0.1 & 80.6$\pm$1.0  \\
    \hline
              & \textbf{Ours(Ent)} & Post-hoc & No & 92.6$\pm$2.0 & 80.8$\pm$3.0 & 81.1$\pm$1.2 &  84.9$\pm$1.2& 85.6$\pm$3.9 \\
              & \textbf{Ours(MaxP)} & Post-hoc & No & 88.6$\pm$3.2 & 78.4$\pm$3.0 & 79.7$\pm$0.6 &  82.4$\pm$0.6& 79.3$\pm$4.5 \\
              & \textbf{Ours(MI)} & Post-hoc & No & 94.3$\pm$1.0 & \textcolor{violet}{\textbf{84.4}$\pm$1.8} & 81.9$\pm$4.7 &  85.5$\pm$3.6 & \textcolor{violet}{\textbf{90.8}$\pm$3.3} \\
              & \textbf{Ours(Dent)} & Post-hoc & No & 93.3$\pm$1.4 & 84.0$\pm$2.3 & 79.5$\pm$3.0 &  84.6$\pm$2.8& 89.8$\pm$2.6 \\
              & \textbf{Ours(Prec)} & Post-hoc & No & \textcolor{violet}{\textbf{94.4}$\pm$1.0} & \textcolor{violet}{\textbf{84.4}$\pm$1.7} & \textcolor{violet}{\textbf{82.1}$\pm$4.9} &  \textcolor{violet}{\textbf{85.6}$\pm$3.6}&
              \textcolor{violet}{\textbf{90.8}$\pm$3.4} \\
    \toprule
  \end{tabular}
  \label{table:OOD}
  \end{center}
  \vspace{-2em}
\end{table*}

\begin{comment}
\begin{table*}[t!]
  \begin{center}
  \scriptsize
  \captionsetup{font=small}
  \caption{\textbf{CIFAR10 OOD Detection AUROC score} }
%   MI, Dent, and Prec stand for different epistemic uncertainty metrics, i.e., Mutual Information, Differential Entropy, and precision. Settings stand for using additional training data/ if training the base model from scratch.
  \vspace{-1.5em}
  \begin{tabular}{lllllllll}
    \\
    \toprule
    \makecell{\textbf{ID Data}\\/\textbf{Model}} & \textbf{Methods} & \textbf{Settings} & \textbf{Additional Data} & \textbf{SVHN}& \textbf{FMNIST} & \textbf{LSUN}& \textbf{TIM} &\textbf{Corrupted} \\
    \hline
    CIFAR10 & Base Model(Ent) & Traditional & No & 86.4$\pm$4.6 & 90.8$\pm$1.3 & 89.0$\pm$0.5 & 87.5$\pm$1.1 & 85.9$\pm$8.2  \\
    VGG16   & Base Model(MaxP) & Traditional & No & 86.3$\pm$4.4 & 90.4$\pm$1.2 & 88.7$\pm$0.5 & 87.3$\pm$1.1 & 85.7$\pm$8.1   \\
              & Whitebox & Post-hoc & Yes & 96.9$\pm$0.9 & 95.2$\pm$1.2 & 89.3$\pm$2.2 & 88.9$\pm$2.5 & 96.4$\pm$1.0   \\
              & LULA & Post-hoc & Yes & 97.1$\pm$1.7 & 94.3$\pm$0.0 & 92.8$\pm$0.1 & 90.0$\pm$0.0 & 97.7$\pm$2.0  \\
    \hline
              & \textbf{Ours(Ent)} & Post-hoc & No & 96.3$\pm$3.0 &89.0$\pm$5.2 & 89.6$\pm$3.4 &  89.4$\pm$3.5& 95.9$\pm$4.3 \\
              & \textbf{Ours(MaxP)} & Post-hoc & No & 95.6$\pm$3.6 & 87.8$\pm$4.4 &89.1$\pm$2.4 &  88.2$\pm$2.6 & 94.0$\pm$7.3 \\
              & \textbf{Ours(MI)} & Post-hoc & No & 100.0$\pm$0.0 & 98.8$\pm$0.5 & 95.2$\pm$0.9 &  98.1$\pm$0.3 & 100.0$\pm$0.0 \\
              & \textbf{Ours(Dent)} & Post-hoc & No & \textcolor{violet}{\textbf{100.0}$\pm$0.0} & \textcolor{violet}{\textbf{98.4}$\pm$0.9} & \textcolor{violet}{\textbf{95.7}$\pm$0.8} &  \textcolor{violet}{\textbf{97.7}$\pm$0.5} & \textcolor{violet}{\textbf{100.0}$\pm$0.0} \\
               & \textbf{Ours(Prec)} & Post-hoc & No & 100.0$\pm$0.0 & 98.8$\pm$0.5 & 95.1$\pm$0.5 &  98.1$\pm$0.3 & 100.0$\pm$0.0 \\
    \toprule
  \end{tabular}
  \label{table:OOD-CIFAR10}
  \end{center}
  \vspace{-1.5em}
\end{table*}

\begin{table*}[t!]
  \begin{center}
  \scriptsize
  \captionsetup{font=small}
  \caption{\textbf{CIFAR100 OOD Detection AUROC score}}
  \vspace{-1.5em}
  \begin{tabular}{lllllllll}
    \\
    \toprule
    \makecell{\textbf{ID Data}\\/\textbf{Model}} & \textbf{Methods} & \textbf{Settings} & \textbf{Additional Data}  & \textbf{SVHN}& \textbf{FMNIST} & \textbf{LSUN}& \textbf{TIM} &\textbf{Corrupted} \\
    \hline
    CIFAR100 & Base Model(Ent) & Traditional & No & 76.2$\pm$5.2 & 77.8$\pm$2.4 & 80.1$\pm$0.5 & 79.7$\pm$0.3 & 65.8$\pm$11.4  \\
    WideResNet   & Base Model(MaxP) & Traditional & No & 73.9$\pm$4.3 & 76.4$\pm$2.3 & 78.7$\pm$0.5 & 78.0$\pm$0.2 & 63.8$\pm$10.4   \\
              & Whitebox & Post-hoc & Yes & 89.0$\pm$0.7 & 82.4$\pm$1.1 & 80.5$\pm$0.7 & 79.0$\pm$1.1 & 83.1$\pm$1.6   \\
              & LULA & Post-hoc & Yes & 84.2$\pm$1.0 & 83.2$\pm$0.1 & 79.6$\pm$0.3 & 78.5$\pm$0.1 & 80.6$\pm$1.0  \\
    \hline
              & \textbf{Ours(Ent)} & Post-hoc & No & 92.6$\pm$2.0 & 80.8$\pm$3.0 & 81.1$\pm$1.2 &  84.9$\pm$1.2& 85.6$\pm$3.9 \\
              & \textbf{Ours(MaxP)} & Post-hoc & No & 88.6$\pm$3.2 & 78.4$\pm$3.0 & 79.7$\pm$0.6 &  82.4$\pm$0.6& 79.3$\pm$4.5 \\
              & \textbf{Ours(MI)} & Post-hoc & No & 94.3$\pm$1.0 & \textcolor{violet}{\textbf{84.4}$\pm$1.8} & 81.9$\pm$4.7 &  85.5$\pm$3.6 & \textcolor{violet}{\textbf{90.8}$\pm$3.3} \\
              & \textbf{Ours(Dent)} & Post-hoc & No & 93.3$\pm$1.4 & 84.0$\pm$2.3 & 79.5$\pm$3.0 &  84.6$\pm$2.8& 89.8$\pm$2.6 \\
              & \textbf{Ours(Prec)} & Post-hoc & No & \textcolor{violet}{\textbf{94.4}$\pm$1.0} & \textcolor{violet}{\textbf{84.4}$\pm$1.7} & \textcolor{violet}{\textbf{82.1}$\pm$4.9} &  \textcolor{violet}{\textbf{85.6}$\pm$3.6}&
              \textcolor{violet}{\textbf{90.8}$\pm$3.4} \\
    \toprule
  \end{tabular}
  \label{table:OOD-CIFAR100}
  \end{center}
  \vspace{-1.5em}
\end{table*}

\end{comment}

\subsection{Settings}\label{sec:setting}
\paragraph{Benchmark.}
For both OOD detection and misclassification detection tasks, we employ three standard datasets to train the base model and the meta-model: MNIST, CIFAR10, and CIFAR100. For each dataset, we use different base-model structures, i.e., LeNet for MNIST, VGG-16~\citep{simonyan2014very} for CIFAR10, and WideResNet-16~\citep{zagoruyko2016wide} for CIFAR100. For LeNet and VGG-16, the meta-model uses extracted feature after each pooling layer, and for WideResNet-16, the meta-model uses extracted feature after each residual black. In general, the total number of intermediate features is less than 5 to ensure computational efficiency. For the OOD task, we consider five different OOD datasets for evaluating the OOD detection performance: Omiglot, FashionMNIST, KMNIST, CIFAR10, and corrupted MNIST as outliers for the MNIST dataset; SVHN, FashionMNIST, LSUN, TinyImageNet, and corrupted CIFAR10 (CIFAR100) as outliers for CIFAR10 (CIFAR100) dataset. For the trustworthy transfer learning task, we use the ResNet-50 pretrained on ImageNet as our pretrained source domain model and adapt the source model to the two target tasks, STL10 and CIFAR10, by training the meta-model.

\paragraph{Baselines.}
%For OOD and misclassification tasks, we compare our proposed method with the following baseline methods: (1) the naive base-model trained with cross-entropy loss; (2) The standard Bayesian method Monte-Carlo Dropout~\citep{gal2016dropout}; (3) The Dirichlet network with variational inference (Be-Bayesian)~\citep{joo2020being}; (4) The meta-model based method (Whitebox)~\citep{chen2019confidence}; (4) The post-hoc uncertainty quantification using Laplace Approximation (LULA)~\citep{kristiadi2021learnable}. For the trustworthy transfer learning task, since there is no existing work designed for this problem, we compare our method with two simple baselines: (1) Fine-tune the last layer of the source model. (2) Train our proposed meta-model on top of the source model using standard cross-entropy loss.
For OOD and misclassification tasks, except the naive base-model trained with cross-entropy loss, we mainly compare with the existing post-hoc UQ methods as baselines: (1) The meta-model based method (Whitebox)~\citep{chen2019confidence}; (2) The post-hoc uncertainty quantification using Laplace Approximation (LULA)~\citep{kristiadi2021learnable}. In order to further validates our strong empirical performance, we also compare with other SOTA intrinsic UQ methods in the Appendix~\ref{results-appendix}: (1) The standard Bayesian method Monte-Carlo Dropout~\citep{gal2016dropout}; (2) The Dirichlet network with variational inference (Be-Bayesian)~\citep{joo2020being}; (3) The posterior network with density estimation~\citep{charpentier2020posterior}; (4) The robust OOD detection method ALOE~\citep{chen2020robust}.

For the trustworthy transfer learning task, since there is no existing work designed for this problem, we compare our method with two simple baselines: (1) Fine-tune the last layer of the source model. (2) Train our proposed meta-model on top of the source model using standard cross-entropy loss.

\paragraph{Performance Metrics.}
We evaluate the UQ performance by measuring the area under the ROC curve (AUROC) and the area under the Precision-Recall curve (AUPR). The results are averaged over five random trails for each experiment. For the OOD task, we consider the in-distribution test samples as the negative class and the outlier samples as the positive class. For the misclassification task, we consider correctly classified test samples as the negative class and miss-classified test samples as the positive class.

\subsection{OOD Detection}\label{sec:OOD}
The OOD detection results for the three benchmark datasets, MNIST, CIFAR10, and CIFAR100, are shown in Table~\ref{table:OOD}. Additional baseline comparisons are provided in Table~\ref{table:OOD-appendix} in Appendix. Our proposed Dirichlet meta-model method consistently outperforms all the baseline methods in terms of AUROC score (AUPR results are shown in Appendix), including the recent proposed SOTA post-hoc uncertainty learning method LULA. We also evaluate the performance of all the uncertainty metrics defined in section~\ref{sec:metric}, as it can be observed that compared to total uncertainty (Ent and MaxP), epistemic uncertainties (MI, Dent, Prec) can achieve better UQ performance for the OOD detection task. Moreover, our proposed method does not require additional data to train the meta-model. In contrast, Whitebox requires an additional validation set to train the meta-model, and LULA also needs an additional OOD dataset during training to distinguish the in-distribution samples and outliers, which imposes practical limitations.


\subsection{Misclassification Detection}\label{sec:Miss}
The misclassification detection results for the three benchmark datasets, MNIST, CIFAR10, and CIFAR100, are shown in Table~\ref{table:miss-class}. Additional baseline comparisons are provided in Table~\ref{table:miss-class-appendix}. LULA turns out to be a strong baseline for the misclassification detection task. Although our proposed method performs slightly worse than LULA in terms of the AUROC, it outperforms all the baselines in terms of AUPR.

\begin{table*}[t]
  \begin{center}
  \scriptsize
  \captionsetup{font=small}
  \caption{\textbf{Misclassification Results.} Ent and MaxP stand for Entropy and Max Probability, respectively. Settings stand for post-hoc or traditional, i.e., training the entire model from scratch.}
  \vspace{-1.5em}
  \begin{tabular}{llllllllll}
    \\
    \toprule
    \textbf{Methods} & \textbf{Settings} & \textbf{Metric} & \textbf{MNIST}& \makecell{\textbf{CIFAR 10}} & \makecell{\textbf{CIFAR 100}}& \textbf{Metric} &  \textbf{MNIST}& \makecell{\textbf{CIFAR 10}} & \makecell{\textbf{CIFAR 100}}\\
    \hline
    Base Model(Ent)  & Traditional & AUROC & 96.7$\pm$0.9 & 92.1$\pm$0.2 & 87.2$\pm$0.2 & AUPR & 37.4$\pm$4.0 & 47.5$\pm$1.2 & 67.0$\pm$1.2  \\
    Base Model(MaxP) & Traditional & AUROC & 96.7$\pm$0.9 & 92.1$\pm$0.2 & 86.8$\pm$0.2 & AUPR & 39.4$\pm$3.6 & 46.6$\pm$1.1 & 65.7$\pm$1.0  \\
    Whitebox &Post-hoc & AUROC & 94.9$\pm$0.2 & 90.2$\pm$0.1 & 80.3$\pm$0.1 & AUPR & 30.4$\pm$0.3 & 45.7$\pm$0.2 & 52.5$\pm$0.3  \\
    LULA &Post-hoc & AUROC & \textcolor{violet}{\textbf{98.8}$\pm$0.1} & \textcolor{violet}{\textbf{94.5}$\pm$0.0} & \textcolor{violet}{\textbf{87.5}$\pm$0.1} & AUPR & 40.7$\pm$4.2 & 47.3$\pm$0.7 & 66.0$\pm$0.4  \\
    \hline
    \textbf{Ours(Ent)} &Post-hoc & AUROC  & 96.9$\pm$0.6 & 91.1$\pm$0.2 & 83.4$\pm$0.1 & AUPR & 35.6$\pm$4.5 & 50.0$\pm$3.1 & 66.3$\pm$0.4  \\
    \textbf{Ours(MaxP)} &Post-hoc & AUROC  & 97.4$\pm$0.4 & 92.2$\pm$0.7 & 85.8$\pm$0.2 & AUPR & \textcolor{violet}{\textbf{44.5}$\pm$5.1} & \textcolor{violet}{\textbf{54.2}$\pm$3.2} & \textcolor{violet}{\textbf{68.2}$\pm$0.5}  \\
    \toprule
  \end{tabular}
  \label{table:miss-class}
  \end{center}
  \vspace{-1.2em}
\end{table*}

\begin{table*}[t]
  \begin{center}
  \scriptsize
  \captionsetup{font=small}
  \caption{\textbf{Trustworthy Transfer Learning Results.} Ent, MaxP, MI, and Dent stand for different uncertainty measurements, i.e., Entropy, Max Probability, Mutual Information, and Differential Entropy, respectively. We use ResNet-50 pretrained with ImageNet as the source model and FashionMNIST as OOD samples.}
  \vspace{-1.5em}
  \begin{tabular}{lllllllll}
    \\
    \toprule
    \textbf{Methods} &\textbf{Target} &\textbf{Test Acc} & \textbf{AUROC} & \textbf{AUPR}& \textbf{Target} &\textbf{Test Acc}& \textbf{AUROC} &  \textbf{AUPR} \\
    \hline
    FineTune(Ent)  &\textbf{STL10} &48.1$\pm$0.5  & 89.2$\pm$0.9 & 89.3$\pm$1.2 & \textbf{CIFAR10} &65.0$\pm$0.4 &74.8$\pm$1.6 &71.6$\pm$1.8  \\
    FineTune(MaxP) &               &48.1$\pm$0.5  & 81.8$\pm$1.2 & 83.2$\pm$1.7 &                  &65.0$\pm$0.4 &72.7$\pm$1.4&69.4$\pm$1.5  \\
    \hline
    CrossEnt Loss(Ent)  &    &48.0$\pm$0.2  & 88.9$\pm$1.0 & 87.4$\pm$1.5 &  &86.3$\pm$0.1 &85.0$\pm$ 1.0& 82.2$\pm$1.9  \\
    CrossEnt Loss(MaxP) &               &48.0$\pm$0.2  & 84.9$\pm$0.7 & 84.7$\pm$0.8 &                  &86.3$\pm$0.1 &83.1$\pm$0.9& 79.0$\pm$1.6  \\
    %\hline
    %\textbf{Ours(Ent)} &           &48.3 & 91.4 & 91.3 &                   &86.3 &87.8&85.3  \\
    %\textbf{Ours(MaxP)}&           &48.3 & 85.7 & 86.3 &                   &86.3 &84.7&81.3  \\
    % \textbf{Ours(MI)} &          &48.3  & 91.0 & 90.7 &                   &86.3 &89.0&84.5  \\
    %\textbf{Ours(Dent)}&          &48.3  & \textcolor{violet}{\textbf{92.4}} & \textcolor{violet}{\textbf{92.0}} &                   &85.3 &\textcolor{violet}{\textbf{90.5}}&\textcolor{violet}{\textbf{87.3}}  \\
    \hline
    \textbf{Ours(Ent)} &          &47.2$\pm$0.3 & \textcolor{violet}{\textbf{91.8$\pm$0.8}} & \textcolor{violet}{\textbf{91.3$\pm$0.7}} &  &86.6$\pm$0.3 &89.9$\pm$1.3&88.8$\pm$1.5\\
    \textbf{Ours(MaxP)}&          &47.2$\pm$0.3 & 87.3$\pm$1.7 & 87.9$\pm$1.5 &  &86.6$\pm$0.3 &87.6$\pm$1.6&85.5$\pm$2.0  \\
     \textbf{Ours(MI)} &        &47.2$\pm$0.3 & 90.2$\pm$2.0 & 88.7$\pm$2.8 &  &86.6$\pm$0.3 &90.8$\pm$0.9&89.9$\pm$1.2\\
    \textbf{Ours(Dent)}&        &47.2$\pm$0.3 & 91.7$\pm$1.0 & 90.3$\pm$1.7 &  &86.6$\pm$0.3 &\textcolor{violet}{\textbf{92.0$\pm$0.8}}& \textcolor{violet}{\textbf{91.1$\pm$0.8}}\\
    \toprule
  \end{tabular}
  \label{table:transfer}
  \end{center}
  \vspace{-1.5em}
\end{table*}

\subsection{Trustworthy Transfer Learning}\label{sec:transfer}
We use ImageNet pretrained ResNet-50 as our source domain base model and adapt the pretrained model to the target task by training the meta-model using the target domain training data. Unlike traditional transfer learning, which only focuses on testing prediction accuracy on the target task, we also evaluate the UQ ability of the meta-model in terms of OOD detection performance. We use FashionMNIST as OOD samples for both target tasks STL10 and CIFAR10 and evaluate the AUROC score. The results are shown in Table~\ref{table:transfer}. Our proposed meta-model method can achieve comparable prediction performance to the baseline methods and significantly improve the OOD detection performance, which is crucial in trustworthy transfer learning.



\subsection{Discussion}\label{sec:ablation}


In this section, we further investigate our proposed method through an ablation study using the CIFAR10 OOD task. Based on our insights and the empirical results, we conclude the following four key factors in the success of our meta-model based method:

\textbf{Feature Diversity.}
We replace our proposed meta-model structure with a simple linear classifier attached to only the final layer. The ablation results are shown in Table~\ref{table:ablation-structure} as ``\textbf{Linear-Meta}''. It can be observed here and in Figure~\ref{fig: toy_problem} that the performance degrades without using features from all intermediate layers, which further justifies the importance of feature diversity and the effectiveness of our meta-model structure. 

% We replace our proposed meta-model structure with a simple linear classifier attached at the final layer. The ablation results are shown in Table~\ref{table:ablation-structure} as ``\textbf{Last Layer}''. It can be observed the performance degrades without using our designed meta-model, which further justifies the importance of feature diversity and effectiveness of our meta-model structure. This observation can also be interpreted in Figure~\ref{fig: toy_problem}, using multiple feature representation will encourage meta-model parameterize a more sharp Dirichlet distribution for in distribution data.



\textbf{Dirichlet Technique.}
Instead of using a meta-model to parameterize a Dirichlet distribution, we train the meta-model using the standard cross-entropy loss, which simply outputs a categorical label distribution. The ablation results are shown in Table~\ref{table:ablation-structure} as ``\textbf{Cross-Ent}''. It can be shown that performance degrades again because it cannot quantify epistemic uncertainty, which justifies the effectiveness of using Bayesian techniques.

\textbf{Validation for Uncertainty Learning.}
We retrain the last layer of the base model using the cross-entropy loss with the proposed validation trick. The results are shown in Table~\ref{table:ablation-structure} as ``\textbf{LastLayer}''. It turns out even such a naive method can achieve improved performance compared to the base model, which further justifies the effectiveness of the post-hoc uncertainty learning setting, i.e., the benefit of solely focusing on UQ performance at the second stage. This interesting observation inspires us to conjecture that efficiently retraining the classifier of the base model at the second stage will lead to better UQ performance. A theoretical investigation of this observation can be interesting for future work.

% We simply retrain the last layer of base-model using the cross-entropy loss with our proposed validation trick. The results are shown in Table~\ref{table:ablation-structure} as ``\textbf{Cross-Ent}+\textbf{lastLayer}''. It turns out even such naive method can achieve improved performance compare to the base-model, which further justifies the effectiveness of the post-hoc uncertainty learning setting, i.e., the benefit of solely learning uncertainty performance at second-stage without focusing on prediction performance. This interesting observation encourages us to conjecture that efficiently re-training the classifier of base-model at second-stage can be a low-hanging-fruit approach for better uncertainty quantification performance. A theoretical investigation for this observation can be a interesting future work.
\textbf{Data Efficiency.}
Instead of using all the training samples, we randomly choose only $10\%$ samples to train the meta-model. The results are shown in Table~\ref{table:ablation-structure} as ``$10\%$\textbf{data}''. It can be observed that our meta-model requires a small amount of data to achieve comparable performance due to the smaller model complexity. Therefore, our proposed method is also more computationally efficient than the approaches that retrain the whole model from scratch.

\begin{table}[h]
  \begin{center}
  \scriptsize
  \captionsetup{font=small}
  \caption{\textbf{Ablation Study of Meta-model (CIFAR10 AUROC score).} The results are reported as mean over five experiment trails. Error bars are provided in Table~\ref{table:ablation-structure-appendix} in the Appendix.}
  \vspace{-1.5em}
  \begin{tabular}{lllllll}
    \\
    \toprule
    \textbf{Methods} & \textbf{SVHN}& \textbf{FMNIST} & \textbf{LSUN}& \textbf{TIM} &\textbf{Corrupted} \\
    \hline
     Base Model(Ent)  & 86.4 & 90.8 & 89.0 & 87.5 & 85.9  \\
    Base Model(MaxP)  & 86.3 & 90.4 & 88.7 & 87.3 & 85.7   \\
    \hline
            \textbf{Linear-Meta(Ent)}  & 90.4 & 91.3 & 91.5 & 89.5 & 90.4 \\
            \textbf{Linear-Meta(MaxP)}  & 90.1 & 91.5 & 91.4 & 89.6 & 90.1 \\
            \textbf{Linear-Meta(MI)}  & 90.6 & 90.1 & 91.2 & 88.8 & 90.5 \\
            \textbf{Linear-Meta(Dent)}  & 90.4 & 90.7 & 91.4 & 89.2 & 90.3 \\
            \textbf{Linear-Meta(Prec)} & 90.6 & 90.0 & 91.2 & 88.8 & 90.6\\
    \hline
      \textbf{Cross-Ent(Ent)}  & 94.2 & 91.2 & 91.2 &  90.3& 94.7 \\
        \textbf{Cross-Ent(MaxP)}  & 93.3& 91.1 &90.9 &  90.0 & 94.0 \\
    \hline
       \textbf{LastLayer(Ent)}& 93.0& 90.2 & 91.9 &  89.9& 93.1 \\
        \textbf{LastLayer(MaxP)} & 92.9 & 90.5 &91.9&  90.1 & 93.1\\
    \hline        
    \textbf{$10\%$data(Ent)}  & 90.0 &89.1& 88.7 &  88.2& 90.2 \\
              \textbf{$10\%$data(MaxP)} & 90.9 & 88.1&86.9&  86.5& 91.7 \\
              \textbf{$10\%$data(MI)}& 99.9 & 98.1 & 95.4&  97.2 & 99.9 \\
              \textbf{$10\%$data(Dent)}  & 96.7 & 97.4 & 94.3&  95.7 & 96.7\\
               \textbf{$10\%$data(Prec)}  & 99.9& 98.0& 95.4 &  97.3& 99.9 \\
    \hline        
    \textbf{Ours(Ent)}  & 96.3 &89.0& 89.6 &  89.4& 95.9 \\
              \textbf{Ours(MaxP)} & 95.6 & 87.8 &89.1 &  88.2 & 94.0 \\
              \textbf{Ours(MI)}& \textcolor{violet}{\textbf{100.0}} & \textcolor{violet}{\textbf{98.8}} & 95.2 &  \textcolor{violet}{\textbf{98.1}} & \textcolor{violet}{\textbf{100.0}} \\
              \textbf{Ours(Dent)}  & \textcolor{violet}{\textbf{100.0}} & 98.4 & \textcolor{violet}{\textbf{95.7}} &  97.7 & \textcolor{violet}{\textbf{100.0}} \\
            \textbf{Ours(Prec)}  & \textcolor{violet}{\textbf{100.0}} & \textcolor{violet}{\textbf{98.8}} & 95.1 &  \textcolor{violet}{\textbf{98.1}} & \textcolor{violet}{\textbf{100.0}} \\
    \toprule
  \end{tabular}
  \label{table:ablation-structure}
  \end{center}
  \vspace{-1.5em}
\end{table}
